\begin{helper}
For every \(\varepsilon>0\), there are \(\delta>0\) and \(s_0\) such
that, whenever \(s\ge s_0\), \(a\in\mathbb N\) satisfies
\(a\le\delta s\), and \(t\in\{1,\ldots,s-1\}\), there is a real number
\(\lambda\) such that, with \(j=s-t\),
\[
  \binom{s+a}{t}\le 2^{(a+j)\lambda},
\]
\[
  a\lambda\le \frac{\varepsilon s}{4},
\]
and
\[
  j\lambda+\frac{3j}{2}-\frac{j^2}{2}\le \frac{\varepsilon s}{4}.
\]
\end{helper}
\begin{proof}
Apply \(\noderef{F2NearDiagonalSmallXLogBound}\) with \(\varepsilon\).  It
gives \(0<\delta\le1\) such that
\[
  x\log_2(2e/x)\le\varepsilon/4
  \qquad(0<x\le\delta). \tag{1}
\]
Apply \(\noderef{F2NearDiagonalLogSquareBound}\) with \(\varepsilon\).  It
gives \(s_1\) such that, for all \(s\ge s_1\),
\[
  \frac12(\log_2(2es)+3/2)^2\le\varepsilon s/4. \tag{2}
\]
Let \(s_0=\max(s_1,1)\).

Now assume \(s\ge s_0\), \(a\le\delta s\), and
\(1\le t\le s-1\).  Put \(j=s-t\).  Then \(j\ge1\), \(j\le s\), and
\[
  s-j=t.
\]
Define
\[
  \lambda=\log_2\left(\frac{e(s+a)}{a+1}\right).
\]
By \(\noderef{F2NearDiagonalChooseSymmetry}\),
\[
  \binom{s+a}{t}=\binom{s+a}{j+a}.
\]
Since \(j\ge1\), we have \(a+1\le j+a\).  Applying
\(\noderef{F2NearDiagonalBinomialLogBound}\) with
\(n=s+a\), \(m=j+a\), and \(b=a+1\) gives
\[
  \binom{s+a}{j+a}
  \le
  2^{(j+a)\log_2(e(s+a)/(a+1))}
  =
  2^{(a+j)\lambda}.
\]
This proves the required binomial control.

It remains to prove the two logarithmic controls.  Since \(\delta\le1\) and
\(a\le\delta s\), we have \(a\le s\), hence \(s+a\le2s\).  Therefore
\[
  \frac{e(s+a)}{a+1}\le 2es.
\]
Taking base-\(2\) logarithms gives
\[
  \lambda\le \log_2(2es). \tag{3}
\]
Using (3) and \(j\ge0\), \(\noderef{F2NearDiagonalQuadraticMaxBound}\)
gives
\[
  j\lambda+\frac{3j}{2}-\frac{j^2}{2}
  \le \frac12(\log_2(2es)+3/2)^2.
\]
Combining this with (2) gives the desired \(j\)-quadratic bound.

For the \(a\lambda\) bound, first suppose \(a=0\).  Then
\(a\lambda=0\le\varepsilon s/4\).  If \(a>0\), put \(x=a/s\).  Since
\(s\ge1\), \(x>0\), and the hypothesis \(a\le\delta s\) gives
\(x\le\delta\).  Also \(s+a\le2s\) and \(a+1\ge a\), so
\[
  \frac{e(s+a)}{a+1}
  \le \frac{2es}{a}
  =
  \frac{2e}{x}.
\]
Thus \(\lambda\le\log_2(2e/x)\), and by (1),
\[
  x\lambda\le x\log_2(2e/x)\le\varepsilon/4.
\]
Multiplying by \(s\) and using \(sx=a\) yields
\[
  a\lambda\le \varepsilon s/4.
\]
All three required controls hold for the displayed value of \(\lambda\).
\end{proof}
